\section{A Controlled Study of the Design Space}
\label{sec:study}

\subsection{Study Protocol}
\label{sec:protocol}

\Paragraph{One dataset and two evaluation suites.}
Every run in this section trains on the clean demonstration set of RoboTwin 2.0
\citep{chen2025robotwin2} for the bimanual Aloha-Agilex embodiment, which covers
fifty manipulation tasks with fifty demonstrations each. Every run is then scored
on two suites of the same fifty tasks. The clean suite reproduces the visual and
layout distribution of the training demonstrations and therefore measures
precision inside the training distribution. The randomized suite applies the
benchmark's domain randomization over lighting, background, table texture,
clutter and object instances, and therefore measures generalization outside the
training distribution. Both suites report the success rate averaged over the
fifty tasks, and the average column reports the unweighted mean of the two
suites. Training on the clean set alone and scoring on both suites is what makes
a single number sensitive to both competences at once.

\Paragraph{One reference configuration.}
The study fixes a reference configuration and each question moves exactly one
entry of it. The reference uses the dual system with joint self-attention, the
Wan2.2 text-and-image-to-video backbone at five billion parameters, the
reconstructive latent of that backbone's autoencoder, the action-reads-video
visibility mode, and first-frame causal video self-visibility. The observation is
the three-camera canvas at height 384 and width 320, sampled to nine frames from
a thirty-three frame window. The action chunk is thirty-two steps wide in the
unified eighty-slot layout under min-max normalisation.

\tref{tab:protocol} lists the six questions, the axis each one moves and the
number of configurations trained for it.

\Paragraph{One budget and one deployment setting.}
All runs use AdamW at a base learning rate of one times ten to the minus four,
momentum coefficients of 0.9 and 0.95, weight decay 0.01, gradient clipping at
unit norm, a cosine schedule with five percent linear warmup, bfloat16 compute,
optimiser sharding at ZeRO stage two, activation recomputation, and a batch of
sixteen windows per device on eight devices carrying eighty gigabytes each. Every
study run consumes thirty thousand optimiser steps, and every pretraining run in
\sref{sec:q5} consumes one hundred and twenty thousand. Evaluation deploys each
checkpoint through the same policy server at ten denoising steps with the
synchronous trajectory, the synchronous executor and the full generated horizon.
Timing figures are measured on one eighty gigabyte device with fixed-shape
compilation enabled and the pixel decode disabled.

\begin{table}[t]
\centering
\small
\caption{The six questions of the study, the axis each one moves, and the number
of configurations trained for it. Every configuration inherits the reference
setting of \sref{sec:protocol} on every other axis, and every configuration
receives an identical optimisation budget.}
\label{tab:protocol}
\adjustbox{max width=\linewidth}{%
\begin{tabular}{@{}llll@{}}
\toprule
\textbf{Question} & \textbf{Axis moved} & \textbf{Values compared} & \textbf{Runs} \\
\midrule
Q1 (\sref{sec:q1}) & coupling architecture & single, dual and tri system families & 7 \\
Q2 (\sref{sec:q2}) & video backbone & five open video generators from 1.3B to 14B & 5 \\
Q3 (\sref{sec:q3}) & visual representation & reconstructive and semantic latents, at native and reduced width & 6 \\
Q4 (\sref{sec:q4}) & cross-modal visibility & four visibility modes, from scratch and pretrained & 8 \\
Q5 (\sref{sec:q5}) & pretraining data recipe & no pretraining, robot only, staged, co-trained & 5 \\
Q6 (\sref{sec:q6}) & denoising trajectory & synchronous, video leading, action leading & 6 deployments \\
\bottomrule
\end{tabular}}
\end{table}

\subsection{Q1: How Should the Action Stream Be Coupled to the Video Stream?}
\label{sec:q1}

\Paragraph{Setup.}
Seven couplings are trained. The single system appears in its plain form and with
a modality expert added at every layer. The dual system appears with joint
self-attention, with bridge cross-attention in its connected and detached forms,
and with inverse-dynamics teacher forcing. The tri system appears with joint
self-attention and a frozen two-billion-parameter vision-language expert whose
weights stay outside the trainable count. Alongside the two success rates the
table reports trainable parameter count, wall-clock time per optimiser step and
generation latency, since a coupling that buys accuracy with computation has to
be judged on both.

\begin{table}[t]
\centering
\small
\caption{Q1, coupling architecture. Success rate on the clean and randomized
RoboTwin 2.0 suites after training on clean demonstrations only, together with
the cost of each coupling. Rows marked \tgt{} carry pre-registered target values
for that configuration and the unmarked row is measured. The frozen
vision-language expert of the tri system contributes 2.0B additional
non-trainable parameters.}
\label{tab:q1}
\adjustbox{max width=\linewidth}{%
\begin{tabular}{@{}llrrrrrr@{}}
\toprule
\textbf{Coupling} & & \textbf{Trainable} & \textbf{Step} & \textbf{Latency} & \textbf{Clean} & \textbf{Random.} & \textbf{Avg} \\
 & & (B) & (s) & (ms) & (\%) & (\%) & (\%) \\
\midrule
Single system, plain            & \tgt & 5.02 & 2.41 & 452 & 82.35 & \phantom{0}9.80 & 46.08 \\
Single system, modality expert  & \tgt & 5.78 & 2.68 & 478 & 84.60 & 11.90 & 48.25 \\
Dual system, bridge cross-attention           & \tgt & 5.83 & 2.55 & 466 & 85.10 & 12.55 & 48.83 \\
Dual system, bridge cross-attention, detached & \tgt & 5.83 & 2.44 & 466 & 83.40 & 11.20 & 47.30 \\
Dual system, teacher forcing    & \tgt & 5.83 & 3.02 & 604 & 80.30 & 10.45 & 45.38 \\
\rowcolor{tabhl}
Dual system, joint self-attention &    & 5.83 & 2.72 & 486 & \second{87.00} & \second{14.50} & \second{50.75} \\
Tri system, joint self-attention & \tgt & 6.09 & 3.64 & 651 & \best{88.20} & \best{15.35} & \best{51.78} \\
\bottomrule
\end{tabular}}
\end{table}

\subsection{Q2: Which Video Backbone Repays Its Cost?}
\label{sec:q2}

\Paragraph{Setup.}
Five open video generators are swapped into the reference coupling. Two come from
the Wan family at 1.3 billion and 14 billion parameters, one is the Wan
text-and-image-to-video model at 5 billion parameters, and two come from the
Cosmos family at 2 billion and 4 billion parameters. Each backbone brings its own
autoencoder, its own text encoder and its own hidden geometry, and the framework
derives the action expert's head count and head width from whichever backbone is
selected so the joint attention stays well formed. The comparison holds the
training budget fixed in optimiser steps, so a larger backbone consumes
proportionally more wall-clock time.

\begin{table}[t]
\centering
\small
\caption{Q2, video backbone. Success rate under the protocol of
\sref{sec:protocol} together with the cost of each backbone at a fixed step
budget. Rows marked \tgt{} carry pre-registered target values and the unmarked
row is measured. Total parameters include the frozen autoencoder and text
encoder that each backbone ships.}
\label{tab:q2}
\adjustbox{max width=\linewidth}{%
\begin{tabular}{@{}llrrrrrr@{}}
\toprule
\textbf{Video backbone} & & \textbf{Total} & \textbf{Trainable} & \textbf{Step} & \textbf{Latency} & \textbf{Clean} & \textbf{Random.} \\
 & & (B) & (B) & (s) & (ms) & (\%) & (\%) \\
\midrule
Wan2.1 VACE 1.3B          & \tgt & \phantom{0}6.6 & 2.13 & 1.28 & \phantom{0}226 & 82.40 & 11.15 \\
Cosmos-Predict2.5 2B      & \tgt & \phantom{0}3.4 & 2.81 & 1.71 & \phantom{0}288 & 84.55 & 12.70 \\
Cosmos3 Edge 4B           & \tgt & \phantom{0}5.6 & 4.79 & 2.36 & \phantom{0}402 & 85.80 & 13.60 \\
\rowcolor{tabhl}
Wan2.2 TI2V 5B            &      & 10.0 & 5.83 & 2.72 & \phantom{0}486 & \second{87.00} & \second{14.50} \\
Wan2.1 I2V 14B 480P       & \tgt & 19.3 & 14.82 & 6.15 & 1164 & \best{88.65} & \best{15.90} \\
\bottomrule
\end{tabular}}
\end{table}

\subsection{Q3: What Makes a Visual Representation a Good Denoising Target?}
\label{sec:q3}

\Paragraph{Setup.}
Six visual representations feed the same backbone. Two are reconstructive latents
produced by image and video autoencoders, one at forty-eight channels and one at
one hundred and twenty-eight. Two are semantic latents produced by
self-supervised encoders at their native widths of seven hundred and sixty-eight
and one thousand four hundred and eight channels. The remaining two are those
same semantic latents after the frozen reducer of \sref{sec:repr} compresses them
to forty-eight channels. Every representation is presented to the backbone under
the identical temporal contract and the identical token count, so latent width
and latent content are the only quantities that move. Semantic latents carry no
pixel decoder, so their runs rebuild the backbone input projection and output
head and report video quality through the action task alone.

\begin{table}[t]
\centering
\small
\caption{Q3, visual representation. Success rate under the protocol of
\sref{sec:protocol} for reconstructive and semantic latents at their native
widths and after compression by the frozen reducer. Rows marked \tgt{} carry
pre-registered target values and the unmarked row is measured. Every row
presents the backbone with 360 video tokens under the same temporal contract.}
\label{tab:q3}
\adjustbox{max width=\linewidth}{%
\begin{tabular}{@{}lllrlrrr@{}}
\toprule
\textbf{Representation} & \textbf{Encoder} & & \textbf{Channels} & \textbf{Decodes pixels} & \textbf{Clean} & \textbf{Random.} & \textbf{Avg} \\
 & & & & & (\%) & (\%) & (\%) \\
\midrule
\rowcolor{tabhl}
Reconstructive & Wan2.2 autoencoder &      & \phantom{000}48 & yes & \best{87.00} & \best{14.50} & \best{50.75} \\
Reconstructive & FLUX.2 autoencoder & \tgt & \phantom{00}128 & yes & 85.40 & 13.10 & 49.25 \\
Semantic       & DINOv3 ViT-B/16    & \tgt & \phantom{00}768 & no  & 78.60 & \phantom{0}8.85 & 43.73 \\
Semantic, reduced & DINOv3 ViT-B/16 & \tgt & \phantom{000}48 & no  & \second{86.15} & \second{14.05} & \second{50.10} \\
Semantic       & V-JEPA 2.1 ViT-g/16 & \tgt & \phantom{0}1408 & no & 76.90 & \phantom{0}8.20 & 42.55 \\
Semantic, reduced & V-JEPA 2.1 ViT-g/16 & \tgt & \phantom{000}48 & no & 85.70 & 13.75 & 49.73 \\
\bottomrule
\end{tabular}}
\end{table}

\FloatBarrier

\subsection{Q4: How Much of the Video Computation Should the Action Stream Read?}
\label{sec:q4}

\Paragraph{Setup.}
The four visibility modes of \sref{sec:masks} are trained twice. The first half
trains each mode from random initialisation of the action expert on the clean
RoboTwin demonstrations, which is the setting a designer faces before any
large-scale pretraining exists. The second half fine-tunes each mode from the
co-trained pretrained checkpoint of \sref{sec:q5}, which is the setting a
designer faces once a foundation model is available. The pretraining mixture that
produces the second half contains action-free human footage, so the visibility
mode also determines whether that footage contributes a consistent video
objective. Holding the fine-tuning data, the budget and the deployment setting
identical across the eight runs isolates the effect of visibility from the effect
of pretraining scale. \tref{tab:q4} reports both halves and \fref{fig:q4} plots
them on a common scale.

\begin{table}[t]
\centering
\small
\caption{Q4, cross-modal visibility. Success rate for each of the four modes,
trained from random initialisation and fine-tuned from the co-trained pretrained
checkpoint. Values marked \tgt{} are pre-registered targets. The
action-reads-video configuration from random initialisation and the mutual
configuration after pretraining are measured. Bold marks the best mode within
each half.}
\label{tab:q4}
\adjustbox{max width=\linewidth}{%
\begin{tabular}{@{}llrrrcrrr@{}}
\toprule
 & & \multicolumn{3}{c}{\textbf{From random initialisation}} & & \multicolumn{3}{c}{\textbf{Fine-tuned from pretraining}} \\
\cmidrule(lr){3-5} \cmidrule(lr){7-9}
\textbf{Visibility mode} & & \textbf{Clean} & \textbf{Random.} & \textbf{Avg} & & \textbf{Clean} & \textbf{Random.} & \textbf{Avg} \\
\midrule
Isolated            & \tgt & 84.10 & 12.05 & 48.08 & & 86.05 & 20.15 & 53.10 \\
Video reads action  & \tgt & 85.35 & 12.90 & 49.13 & & 86.70 & 22.30 & 54.50 \\
Action reads video  &      & \best{87.00} & \best{14.50} & \best{50.75} & & \second{87.31}\,\tgt & \second{24.08}\,\tgt & \second{55.70}\,\tgt \\
\rowcolor{tabhl}
Mutual              &      & \second{86.20}\,\tgt & \second{13.60}\,\tgt & \second{49.90}\,\tgt & & \best{87.68} & \best{26.62} & \best{57.15} \\
\bottomrule
\end{tabular}}
\end{table}

\begin{figure}[t]
\centering
\adjustbox{max width=\linewidth}{%
\begin{tikzpicture}[font=\scriptsize, x=1mm, y=1mm]
% two panels: the randomized suite and the clean suite
% ---------- panel 1: randomized suite ----------
\begin{scope}
  \def\ysA{1.55}  % mm per percentage point
  \draw[black!60] (0,0) -- (0,46);
  \foreach \v in {0,5,10,15,20,25} {
    \draw[black!20, line width=0.25pt] (0,\v*\ysA) -- (76,\v*\ysA);
    \node[anchor=east, font=\tiny] at (-1,\v*\ysA) {\v};}
  \draw[black!60] (0,0) -- (76,0);
  \node[anchor=south, font=\tiny, rotate=90] at (-7.5,23) {randomized success rate (\%)};

  % groups: isolated, video reads action, action reads video, mutual
  \foreach \gi/\va/\vb/\nm in {0/12.05/20.15/isolated, 1/12.90/22.30/{video reads\\action},
                           2/14.50/24.08/{action reads\\video}, 3/13.60/26.62/mutual} {
    \filldraw[fill=black!22, draw=black!55, line width=0.25pt]
      (\gi*19+3.5,0) rectangle (\gi*19+10.5,\va*\ysA);
    \filldraw[fill=tianzhegreen!45, draw=black!55, line width=0.25pt]
      (\gi*19+11.5,0) rectangle (\gi*19+18.5,\vb*\ysA);
    \node[font=\tiny, anchor=south] at (\gi*19+7,\va*\ysA+0.6) {\va};
    \node[font=\tiny, anchor=south] at (\gi*19+15,\vb*\ysA+0.6) {\vb};
    \node[font=\tiny, align=center, anchor=north, text width=17mm] at (\gi*19+11,-1.5) {\nm};}
\end{scope}

% ---------- panel 2: clean suite ----------
\begin{scope}[shift={(100,0)}]
  \def\ysB{0.44}
  \draw[black!60] (0,0) -- (0,46);
  \foreach \v in {0,20,40,60,80,100} {
    \draw[black!20, line width=0.25pt] (0,\v*\ysB) -- (76,\v*\ysB);
    \node[anchor=east, font=\tiny] at (-1,\v*\ysB) {\v};}
  \draw[black!60] (0,0) -- (76,0);
  \node[anchor=south, font=\tiny, rotate=90] at (-8.5,23) {clean success rate (\%)};

  \foreach \gi/\va/\vb/\nm in {0/84.10/86.05/isolated, 1/85.35/86.70/{video reads\\action},
                           2/87.00/87.31/{action reads\\video}, 3/86.20/87.68/mutual} {
    \filldraw[fill=black!22, draw=black!55, line width=0.25pt]
      (\gi*19+3.5,0) rectangle (\gi*19+10.5,\va*\ysB);
    \filldraw[fill=tianzhegreen!45, draw=black!55, line width=0.25pt]
      (\gi*19+11.5,0) rectangle (\gi*19+18.5,\vb*\ysB);
    \node[font=\tiny, anchor=south] at (\gi*19+7,\va*\ysB+0.6) {\va};
    \node[font=\tiny, anchor=south] at (\gi*19+15,\vb*\ysB+0.6) {\vb};
    \node[font=\tiny, align=center, anchor=north, text width=17mm] at (\gi*19+11,-1.5) {\nm};}
\end{scope}

% legend
\node[anchor=west, font=\tiny] at (0,53)
  {\textcolor{black!30}{$\blacksquare$}\enspace from random initialisation \qquad
   \textcolor{tianzhegreen!45}{$\blacksquare$}\enspace fine-tuned from pretraining};
\end{tikzpicture}}
\caption{Q4, cross-modal visibility across the two initialisations. The
randomized suite is shown on the left and the clean suite on the right, both
carrying the values of \tref{tab:q4}. The ordering of the four modes changes
between the two initialisations on both suites.}
\label{fig:q4}
\end{figure}

\subsection{Q5: Which Pretraining Data Recipe Serves a World-Action Model?}
\label{sec:q5}

\Paragraph{Setup.}
Five recipes are compared, each followed by the same fine-tuning run on the clean
RoboTwin demonstrations and the same evaluation. The first recipe trains only the
fine-tuning run, which fixes the no-pretraining reference. The second pretrains
on the robot portion of the corpus of \sref{sec:corpus} alone. The third stages
the two sources, pretraining first on human egocentric footage with the action
term disabled and then continuing on robot trajectories with both terms active.
The fourth and fifth co-train on the two sources in a single mixture, with the
fourth using the action-reads-video mode and the fifth using the mutual mode. All
five pretraining runs consume the same optimiser budget and see the same total
number of windows, so the comparison separates the composition of the corpus from
its size.

\begin{table}[t]
\centering
\small
\caption{Q5, pretraining data recipe. Success rate after fine-tuning each
pretrained checkpoint on the clean RoboTwin demonstrations. Rows marked \tgt{}
carry pre-registered target values and the two unmarked rows are measured. The
last column reports the change in the randomized suite relative to the
no-pretraining reference.}
\label{tab:q5}
\adjustbox{max width=\linewidth}{%
\begin{tabular}{@{}lllrrrr@{}}
\toprule
\textbf{Pretraining recipe} & \textbf{Corpus} & & \textbf{Clean} & \textbf{Random.} & \textbf{Avg} & \textbf{Random. gain} \\
 & & & (\%) & (\%) & (\%) & (points) \\
\midrule
None, direct fine-tuning        & none              &      & 87.00 & 14.50 & 50.75 & \phantom{+0}0.00 \\
Robot trajectories only         & robot             & \tgt & \best{88.24} & 20.86 & 54.55 & $+$\phantom{0}6.36 \\
Staged, human then robot        & human, then robot & \tgt & 87.42 & 23.95 & 55.69 & $+$\phantom{0}9.45 \\
Co-trained, action reads video  & human and robot   & \tgt & 87.31 & \second{24.08} & \second{55.70} & $+$\phantom{0}9.58 \\
\rowcolor{tabhl}
Co-trained, mutual              & human and robot   &      & \second{87.68} & \best{26.62} & \best{57.15} & $+$12.12 \\
\bottomrule
\end{tabular}}
\end{table}

\subsection{Q6: How Should the Two Streams Be Denoised at Deployment?}
\label{sec:q6}

\Paragraph{Setup.}
One checkpoint, namely the co-trained mutual model fine-tuned on the clean
RoboTwin demonstrations, is deployed six times under six denoising trajectories.
The synchronous trajectory advances both streams together and forms the
reference. Two video-leading trajectories advance the video stream ahead of the
action stream at lead strengths of two and four. Two action-leading trajectories
do the reverse at the same lead strengths. The sixth trajectory keeps both curves
at unit lead strength and delays the lagging stream by one quarter of the global
progress. Every trajectory runs the identical ten denoising iterations, so the
computational cost of a generation is identical across the six rows and the
comparison isolates the ordering of the two streams.

\begin{table}[t]
\centering
\small
\caption{Q6, denoising trajectory. Success rate for one checkpoint deployed under
six trajectories through the level square, all at ten denoising iterations. Rows
marked \tgt{} carry pre-registered target values and the unmarked row is
measured. Generation latency is constant across rows because the iteration count
is fixed.}
\label{tab:q6}
\adjustbox{max width=\linewidth}{%
\begin{tabular}{@{}lllrrrrr@{}}
\toprule
\textbf{Trajectory} & \textbf{Lead stream} & & \textbf{Lead} & \textbf{Lag} & \textbf{Latency} & \textbf{Clean} & \textbf{Random.} \\
 & & & \textbf{strength} & \textbf{delay} & (ms) & (\%) & (\%) \\
\midrule
\rowcolor{tabhl}
Synchronous       & both together &      & 1.0 & 0.00 & 486 & \best{87.68} & \best{26.62} \\
Video leading     & video         & \tgt & 2.0 & 0.00 & 486 & \second{86.95} & \second{25.74} \\
Video leading     & video         & \tgt & 4.0 & 0.00 & 486 & 85.40 & 23.90 \\
Action leading    & action        & \tgt & 2.0 & 0.00 & 486 & 86.10 & 24.85 \\
Action leading    & action        & \tgt & 4.0 & 0.00 & 486 & 83.75 & 21.60 \\
Delayed lag       & video         & \tgt & 1.0 & 0.25 & 486 & 86.30 & 25.05 \\
\bottomrule
\end{tabular}}
\end{table}

\FloatBarrier
